%------------------------------------------------------------------------------%
\begin{question}
  Dado \(z=\left(\sqrt{3}+i\right)^4 \)

  calcule

  a) a parte real de $z$,

  b) a parte imaginária de $z$,

  c) o módulo de $z$,

  d) o argumento de $z$
\end{question}

%------------------------------------------------------------------------------%
\begin{resolution}{Avaliando $u$ na forma polar}
  Convertendo $u=\sqrt{3}+i$ para a forma polar
  \[
    \pause
    \rho
    \pause
    = \abs{u}
    \pause
    = \sqrt{\left(\sqrt{3}\right)^2 + 1^2}
    \pause
    = \sqrt{4}
    \pause
    = 2
  \]
  \[
    \pause
    \sin(\varphi)
    \pause
    = \frac{\Im(u)}{\abs{u}}
    \pause
    = \frac{1}{2}
    \qquad
    \pause
    \cos(\varphi)
    \pause
    = \frac{\Re(u)}{\abs{u}}
    \pause
    = \frac{\sqrt{3}}{2}
    \qquad
    \pause
    \varphi
    \pause
    = \arg(u)
    \pause
    = \ang{30}
    \pause
    = \frac{\pi}{6}
  \]
  \pause
  Assim
  \[
    u
    = \rho\left[\cos\left(\varphi\right)+ i\sin\left(\varphi\right)\right]
    \pause
    = 2\left[\cos\left(\frac{\pi}{6}\right)+ i\sin\left(\frac{\pi}{6}\right)\right]
  \]
\end{resolution}

%------------------------------------------------------------------------------%
\begin{resolution}{Avaliando $z$}
  \begin{align*}
    \onslide<+->{z & = u^4 \\}
    \onslide<+->{& = \rho^4\left[\cos\left(4\varphi\right)+ i\sin\left(4\varphi\right)\right] \\}
    \onslide<+->{& = 2^4\left[\cos\left(\frac{4\pi}{6}\right)+ i\sin\left(\frac{4\pi}{6}\right)\right] \\}
    \onslide<+->{& = 16\left[\cos\left(\frac{2\pi}{3}\right)+ i\sin\left(\frac{2\pi}{3}\right)\right] \\}
    \onslide<+->{& = 16\left[-\frac{1}{2}+ i\frac{\sqrt{3}}{2}\right] \\}
    \onslide<+->{& = -8 + 8\sqrt{3}i}
  \end{align*}
\end{question}

%------------------------------------------------------------------------------%
\begin{resolution}{Avaliando $z$}
\begin{columns}[t]
\column{0.5\linewidth}
  Parte real de $z$
  \[
    \Re(z) = -8
  \]
  \pause
  Parte imaginária de $z$
  \[
    \Im(z) = 8\sqrt{3}
  \]
\column{0.5\linewidth}
  \pause
  Módulo de $z$
  \[
    \abs{z} = 16
  \]
  \pause
  Argumento de $z$
  \[
    \arg(z) = \frac{2\pi}{3}
  \]
\end{columns}
\end{resolution}

%------------------------------------------------------------------------------%


